\subsubsection{Deducción de las ecuaciones de medición e incertidumbre}

Sea $F: \mathbb{R}^N \to \mathbb{R}$ una función aleatoria continua con densidad de
probabilidad conjunta $P(\vec{x})$ sobre $\Omega$, en el sentido de la
ecuación~\eqref{eq:mean-con-mult}, y sea $\vec{\mu} = (\mu_1, \dots, \mu_N)$ el
vector de valores medios $\mu_j = \bk{x_j}$. Conforme a la
ecuación~\eqref{eq:error-func}, $F$ puede expresarse como la suma de su valor medio
y una variación aleatoria $E(\vec{x})$, de modo que, aplicando la definición de la
ecuación~\eqref{eq:mean-con-mult},

\begin{align}
  F(\vec{x}) &= \bk{F} + E(\vec{x})
  \quad\Longrightarrow\quad
  E(\vec{x}) = F(\vec{x}) - \int_\Omega F(\vec{x})\,P(\vec{x})\,d^Nx ,
  \label{eq:modelo-base}
\end{align}

y de la ecuación~\eqref{eq:error-std} se sabe que $\bk{E} = 0$ y que
$\sigma_F^2 = \bk{E^2}$.

En el proceso de calibración interesa cuantificar el \textbf{error de indicación}.
Identificando la indicación media del pie de rey con $F(\vec{x}) \equiv L_i$ y el
valor de referencia del patrón con $\bk{F} \equiv L_{BP}$, la
ecuación~\eqref{eq:modelo-base} conduce a

\begin{equation}
  L_i = L_{BP} + E \quad\Longrightarrow\quad E = L_i - L_{BP},
  \label{eq:e-base}
\end{equation}

donde $L_i$ es el valor medio indicado por el pie de rey y $L_{BP}$ la longitud del
bloque patrón corregida por certificado. Esta expresión es análoga a la
ecuación~\eqref{eq:error-func} de la sección anterior.

\subsubsection*{Corrección por efectos térmicos}

Sea $T_0 = 20\ ^\circ\text{C}$ la temperatura de referencia y $\Delta T$ la
desviación de la temperatura de ensayo respecto de ella. Si $\alpha_i$ y
$\alpha_{BP}$ son los coeficientes de expansión térmica del pie de rey y del bloque
patrón, respectivamente, ambos cuerpos experimentan una dilatación lineal y la
ecuación~\eqref{eq:e-base} se corrige como

\begin{equation}
  E = L_i (1 + \alpha_i \Delta T_i) - L_{BP} (1 + \alpha_{BP} \Delta T_{BP}),
  \label{eq:e-termal}
\end{equation}

con $\Delta T_i = T_i - T_0$ y $\Delta T_{BP} = T_{BP} - T_0$. Expandiendo el
producto y agrupando términos,

\begin{align}
  E &= L_i + L_i \alpha_i \Delta T_i - L_{BP} - L_{BP} \alpha_{BP} \Delta T_{BP}
  \nonumber\\
  &= L_i - L_{BP} + L_i \alpha_i \Delta T_i - L_{BP} \alpha_{BP} \Delta T_{BP}.
  \label{eq:e-expanded}
\end{align}

Para separar el efecto del material del efecto del gradiente térmico se definen las
dos diferencias auxiliares

\begin{align}
  \delta\alpha &= \alpha_i - \alpha_{BP}, \label{eq:dalpha} \\
  \delta T &= \Delta T_i - \Delta T_{BP}, \label{eq:dT}
\end{align}

de las cuales $\alpha_i = \alpha_{BP} + \delta\alpha$ y
$\Delta T_{BP} = \Delta T_i - \delta T$. Sustituyendo en
\eqref{eq:e-expanded} y desarrollando,

\begin{align}
  E &= L_i + L_{BP}\left(\delta\alpha + \alpha_{BP}\right)\Delta T_i
     - L_{BP} + L_{BP}\alpha_{BP}\Delta T_i
     - L_{BP}\alpha_{BP}\Delta T_i \nonumber\\
  &\quad + L_{BP}\alpha_{BP}\,\delta T
     + \left(L_i - L_{BP}\right)\delta\alpha\,\Delta T_i .
\end{align}

El último término es proporcional al producto del error de indicación por dos
diferencias pequeñas, de modo que es de segundo orden. Despreciándolo y asumiendo la
proximidad dimensional $L_i \approx L_{BP}$ en los términos correctivos, se obtiene
el modelo de medición

\begin{equation}
  \therefore\quad
  E = L_i + L_i \, \delta\alpha \, \Delta T_i + L_{BP} \, \alpha_{BP} \, \delta T
  - L_{BP}.
  \label{eq:e-model}
\end{equation}

\subsubsection*{Relación funcional y cálculo de la incertidumbre combinada}

De la ecuación~\eqref{eq:e-model} se observa que $E$ es una función diferenciable
de seis variables de entrada,

\begin{equation}
  E = f(L_i; L_{BP}; \delta\alpha; \Delta T_i; \alpha_{BP}; \delta T),
  \label{eq:e-func}
\end{equation}

por lo que su varianza se obtiene aplicando la ley de propagación de la
ecuación~\eqref{eq:error-std-values}. Como las variables $x_j$ son estadísticamente
independientes, $\text{Cov}(x_i, x_j) = 0$ para $i \neq j$ y la suma se reduce a

\begin{equation}
  u_c^2(E) = \sum_{j=1}^6 \left. \left( \frac{\partial f}{\partial x_j}
  \right) \right|_{\vec{\mu}}^2 u^2(x_j),
  \label{eq:gum}
\end{equation}

donde los coeficientes de sensibilidad $\pti{f} = c_j$ se evalúan en el vector de
valores medios $\vec{\mu}$ según la ecuación~\eqref{eq:error-std-values}. Desarrollando
la sumatoria~\eqref{eq:gum} sobre las seis variables de entrada del
modelo~\eqref{eq:e-func}:

\begin{align}
  u_c^2(E) &=
  \left( \ptij{E}{L_i} \right)^2 u^2(L_i) +
  \left( \ptij{E}{L_{BP}} \right)^2 u^2(L_{BP}) \nonumber \\
  &\quad + \left( \ptij{E}{\delta\alpha} \right)^2 u^2(\delta\alpha) +
  \left( \ptij{E}{\Delta T_i} \right)^2 u^2(\Delta T_i) \label{eq:gum-expanded} \\
  &\quad + \left( \ptij{E}{\alpha_{BP}} \right)^2 u^2(\alpha_{BP}) +
  \left( \ptij{E}{\delta T} \right)^2 u^2(\delta T). \nonumber
\end{align}

Evaluando las derivadas parciales de la función~\eqref{eq:e-model} respecto a cada
variable de entrada:

\begin{align}
  \ptij{E}{L_i} &= 1 + \delta\alpha \, \Delta T_i, &
  \ptij{E}{L_{BP}} &= -1 + \alpha_{BP} \, \delta T, \nonumber \\
  \ptij{E}{\delta\alpha} &= L_i \, \Delta T_i, &
  \ptij{E}{\Delta T_i} &= L_i \, \delta\alpha, \label{eq:sens-deriv} \\
  \ptij{E}{\alpha_{BP}} &= L_{BP} \, \delta T, &
  \ptij{E}{\delta T} &= L_{BP} \, \alpha_{BP}. \nonumber
\end{align}

Sustituyendo los coeficientes~\eqref{eq:sens-deriv} en la
ecuación~\eqref{eq:gum-expanded} se obtiene la ecuación final para la varianza
combinada del error de indicación:

\begin{align}
  u_c^2(E) &=
  [1 + \delta\alpha \, \Delta T_i]^2 u^2(L_i) +
  [-1 + \alpha_{BP} \, \delta T]^2 u^2(L_{BP}) \nonumber \\
  &\quad + [L_i \, \Delta T_i]^2 u^2(\delta\alpha) +
  [L_i \, \delta\alpha]^2 u^2(\Delta T_i) \label{eq:sigma-final} \\
  &\quad + [L_{BP} \, \delta T]^2 u^2(\alpha_{BP}) +
  [L_{BP} \, \alpha_{BP}]^2 u^2(\delta T). \nonumber
\end{align}
